\begin{abstract}

本模板是面向中国原子能科学研究院（CIAE）硕士、博士研究生的学位论文 \LaTeX{} 排版模板，
基于中国科学技术大学论文模板（ustcthesis）修改而来，
在保留原模板完善排版功能的基础上，针对原子能院的封面格式与论文规范进行了适配。

本模板的目标是将论文格式与论文内容分离，
使作者专注于论文内容的写作，而无需过多关注排版细节。
模板已预设符合学院要求的页面尺寸、字体字号、页眉页脚、参考文献格式等，
开箱即用，按照各章注释提示填写内容即可。

% -------------------------------------------------------
% 摘要写作说明：
%   1. 中文摘要一般为 500 字左右（硕士），博士可适当增加。
%   2. 摘要应概括研究目的、方法、主要结果和结论，不宜涉及推导过程。
%   3. 避免出现参考文献、图表编号等。
%   4. 关键词 3～8 个，各词之间用 、 分隔。
% -------------------------------------------------------

\keywords{中国原子能科学研究院、学位论文、 \LaTeX{}、模板、排版}
\end{abstract}

\begin{enabstract}

This template is a \LaTeX{} dissertation typesetting template for master's and doctoral
students at the China Institute of Atomic Energy (CIAE).
It is adapted from the ustcthesis template of the University of Science and Technology of China,
with modifications to cover formats and thesis specifications tailored to CIAE requirements.

The goal of this template is to separate document formatting from content,
allowing authors to focus on writing without worrying about typesetting details.
Page layout, fonts, headers and footers, and reference styles are all preset
to meet institutional requirements.

% -------------------------------------------------------
% English abstract notes:
%   1. The English abstract should match the Chinese version in content.
%   2. Aim for around 300 words for a master's thesis.
%   3. Use the past tense when describing methods and results.
%   4. Keywords should correspond to the Chinese keywords.
% -------------------------------------------------------

\enkeywords{China Institute of Atomic Energy, Dissertation, \LaTeX{} Template, Typesetting}
\end{enabstract}
